\subsection{Fixed MSA adapters on a retrospective CASP16 subset}
\label{app:casp16}

We evaluated the existing Protenix Mini-default v0.5.0 MSA32 and MSA256
Factor/G+ checkpoints, with three training seeds per head and cap and no
retraining. Matched native baselines used caps of 32 and 256 or full MSA.
All conditions used three paired inference seeds, four recycle passes,
five diffusion steps, one structure per seed and target, FP32 and no templates.
The 15 model/condition entries gave 1,125 predictions over 25 targets.

\paragraph{Selection and scope.}
We screened the existing official CASP16.T1 sequence catalog for lengths
128--384, standard amino acids, exact reference sequence/residue-number
agreement and at least 80\% C$\alpha$ reference coverage. A global BLOSUM62
alignment with gap-open $-10$ and extension $-0.5$ excluded matches to
Train384 or Dev8 with at least 30\% identity and 80\% coverage of both
sequences. Of 31 length/alphabet-eligible targets, four failed reference
coverage and two failed the homology screen. Selection used no model scores.
The primary subset contained eight official A1 targets without component
suffixes: T1207, T1231, T1246, T1266, T1276, T1278, T1280 and T1299.
Only the target protein, MSA and adapter features were supplied, including
when the experimental context contained DNA or ligands. The other 17
targets comprised 12 complex components and five homooligomer/other chains;
their scores remain descriptive and are not pooled into the primary subset.
This retrospective, length-restricted evaluation neither represents all
CASP16 targets nor isolates foundation-model pretraining or historical MSA
database exposure. It is not a blind CASP benchmark.

\paragraph{Endpoints and uncertainty.}
The two co-primary endpoints were Factor-minus-matched-baseline C$\alpha$
pair-lDDT at the two caps. After averaging training and paired inference
seeds within each of the eight targets, 20,000 paired target-bootstrap
draws provided Bonferroni 97.5\% intervals for these two gains. Intervals
condition on the fixed models and inference seeds; eight targets limit
precision. Other contrasts and metrics use unadjusted 95\% intervals.
MSA assembly retained the original query row and at most 31 or 255 source
sequences; recorded actual depths, rather than caps, describe each input.
Baseline and adapter conditions shared the same processed MSA inputs.

\begin{table}[!htbp]\centering\small
\setlength{\tabcolsep}{4pt}
\caption{CASP16 primary subset ($N=8$), C$\alpha$ pair-lDDT. $B$: matched
unadapted Mini-default; $F/G^+$: fixed Native heads. The two $F-B$ intervals
are Bonferroni 97.5\%; $F-G^+$ intervals are unadjusted 95\%. Full MSA is
a descriptive native reference, with no full-MSA adapter evaluated here.}
\label{tab:casp16_pair}
\begin{tabular}{@{}lrrrrr@{}}\toprule
MSA input & $B$ & $F$ & $G^+$ & $F-B$ & $F-G^+$ \\\midrule
\rebuttalCaspPairRows
\bottomrule\end{tabular}\end{table}

The 32-row Factor gain passes the prespecified correction; the 256-row
gain does not. G+ also gains at 32 rows under its unadjusted secondary
interval. Neither cap establishes Factor superiority over G+, and the
full-MSA baseline has the highest mean. These findings support capped-MSA
adaptation on this subset without establishing a unique Factor benefit.

\begin{table}[!htbp]\centering\small
\setlength{\tabcolsep}{4pt}
\caption{Secondary CASP16 metrics on the same eight primary targets.
All contrast intervals are unadjusted 95\% target-bootstrap intervals.}
\label{tab:casp16_other}
\begin{tabular}{@{}lrrrrr@{}}\toprule
MSA input & $B$ & $F$ & $G^+$ & $F-B$ & $F-G^+$ \\\midrule
\multicolumn{6}{@{}l}{\emph{Residue-average C$\alpha$ lDDT}} \\
\rebuttalCaspResidueRows
\midrule
\multicolumn{6}{@{}l}{\emph{TM-score, normalized by full target length}} \\
\rebuttalCaspTmRows
\bottomrule\end{tabular}\end{table}
